\subsection{有限性判别准则}

我们沿用前面各号的记号，并假设 S 有限。

定理 2. 以下性质等价：

\begin{enumerate}
\def\labelenumi{(\arabic{enumi})}
\item
  W 是有限的。
\item
  型 \(B_{M}\) 是正的且非退化的。
\item
  \(\Longrightarrow\) (2). 设 \(S = \bigcup_{i} S_{i}\) 是 \(S\) 分解为连通分支的分解（第 IV 章，§ 1，第 9 号），并设 \(W = \prod_{i} W_{i}\) 是 \(W\) 的相应分解。空间 \(E\) 可与诸空间 \(E_{i} = R^{S_{i}}\) 的直和等同，而 \(B_{M}\) 可与诸相应的型 \(B_{M_{i}}\) 的直和等同。于是我们化归到 \((W, S)\) 不可约的情形。由于假设 \(W\) 有限，\(E\) 是半单的 \(W\) -模（附录，命题 2）。由命题 5 的推论可知 \(E\) 是绝对单的。设 \(B'\) 是 \(E\) 上一个正的非退化型，并设 \(B''\) 是它在 \(W\) 下诸变换的和。由于 \(B''\) 在 \(W\) 下不变，它与 \(B_{M}\) 成比例（ \(\S 2\) ，第 1 号，命题 1）；由于 \(B_{M}(e_{s}, e_{s}) = 1\) （对所有 \(s \in S\) ）成立，比例系数 \(>0\) ，而由于 \(B''\) 是正的，\(B_{M}\) 亦然，这就证明了 (2)。
\item
  \(\Longrightarrow\) (1). 若 \(B_{M}\) 正且非退化，则正交群 \(\mathbf{O}(B_{M})\) 是紧的（《积分》第 VII 章，§3，第 1 号）。由于 \(\sigma(W)\) 是 \(\mathbf{O}(B_{M})\) 的离散子群（定理 1 的推论 3），可知 \(\sigma(W)\) 有限，从而 W 亦有限。证毕。
\end{enumerate}

下面的结果已在证明过程中得到建立：

推论. 若 \((\mathrm{W},\mathrm{S})\) 不可约且有限，则 \(\mathbf{E}\) 是绝对单的 W-模。

定理 2 所提供的有限性判别准则使得可以对所有有限 Coxeter 群进行分类（参见~第 VI 章，§4）。我们在这里只限于下面的初步结果：

命题 8. 若 W 有限，则 (W, S) 的图是一个森林（第 IV 章，附录）。

若不然，这个图将包含一个回路 \((s_{1},\ldots,s_{n})\) ，其中 \(n\geqslant3\) 。置 \(m_{i}=m(s_{1},s_{i+1})\) （\(1\leqslant i<n\) ）及 \(m_{n}=m(s_{n},s_{1})\) ，这就意味着对所有 i 有 \(m_{i}\geqslant3\) 。设

\[
x = e _ {s _ {1}} + \dots + e _ {s _ {n}}.
\]

于是 \(\mathrm{B_M}(x,x) = n + 2\sum_{i < j}\mathrm{B_M}(e_{s_i},e_{s_j})\) 。而

\[
\mathrm{B} _ {\mathrm{M}} (e _ {s _ {i}}, e _ {s _ {i + 1}}) = - \cos \frac {\pi}{m _ {i}} \leqslant - \cos \frac {\pi}{3} \leqslant - \frac {1}{2},
\]

对 \(\mathrm{B}_{\mathbb{M}}(e_{s_n}, e_{s_i})\) 同样如此。由于和中其余各项 \(\leqslant 0\) ，我们得到

\[
\mathrm{B} _ {\mathrm{M}} (x, x) \leqslant n - n = 0,
\]

这与 \(B_{M}\) 正且非退化相矛盾。

推论. 若 (W, S) 不可约且有限，则其图是一棵树。

事实上，连通的森林是树。

与 §3 结果的比较。

首先设 \((W,S)\) 是一个有限 Coxeter 群。用 \((x|y)\) 表示型 \(\mathrm{B}_{\mathbf{M}}(x,y)\) ；由定理 2，这是 E 上的一个标量积。对所有 \(s \in S\) ，设 \(H_{s}\) 是与正交反射 \(\sigma_{s}\) 相伴随的超平面，并设 \(\mathfrak{H}\) 为超平面族 \(w(\mathrm{H}_{s})\) （其中 \(s \in S\) ，\(w \in W\) ）。设 \(C_{0}\) 为由 \(x \in E\) （满足 \((x|e_{s}) > 0\) ，对所有 \(s \in S\) ）所成的集合。最后，借助 \(\sigma\) 把 W 与空间 E 的正交群 \(\mathbf{O}(\mathbf{E})\) 的一个子群等同起来。

命题 9. 按照前面的记号，W 是 O(E) 的由关于 \(\mathfrak{H}\) 中诸超平面的反射所生成的子群。它是一个本质的群（§3，第 7 号），且 C0 是 E 相对于 \(\mathfrak{H}\) 的一个室。

第一个断言是平凡的。另一方面，若 \(x \in E\) 在 W 下不变，则它与所有 \(e_{s}\) 正交，从而为零；这表明 W 是本质的。最后，同构 \(E \to E^{*}\) （由 \(B_{M}\) 定义）把 \(C_{0}\) 变为第 4 号中的集合 C；在那里证明的性质 \((P_{n})\) 表明，对所有 \(w \in W\) 与所有 \(s \in S\) ，\(w(C_0)\) 不与 \(H_s\) 相交。由此可知 \(C_0\) 含于 \(\mathfrak{H}\) 中诸超平面之并的补集 U，而由于 \(C_0\) 在 U 中连通、开且闭，它是 E 相对于 \(\mathfrak{H}\) 的一个室。证毕。

我们可以把 §3 中证明的所有性质应用于 W 与 \(C_{0}\) 。特别地，\(\overline{C_{0}}\) 是 W 在 E 上作用的一个基本域（换句话说，第 6 号中定义的锥 U 等于整个 E）。

反之，设 E 是一个有限维实向量空间，配备标量积 \((x|y)\) ，并设 W 是 E 的一个本质的有限位移群且保持 0 固定；假设 W 由反射生成。设 \(C_{0}\) 是 E 相对于 W 的一个室（参见~§3），并设 S 是相对于 \(C_{0}\) 的墙的正交反射所成的集合。则 (W, S) 是一个有限 Coxeter 系统（§3，第 2 号，定理 1）。此外，若 \(s \in S\) ，用 \(H_{s}\) 表示 \(C_{0}\) 的与 s 相应的墙，并用 \(e_{s}\) 表示正交于 \(H_{s}\) 且关于 \(H_{s}\) 与 \(C_{0}\) 位于同侧的单位向量。若 \((m(s, s'))\) 表示 (W, S) 的 Coxeter 矩阵，则 §3 的命题 3 与 7 表明

\[
(e _ {s} | e _ {s ^ {\prime}}) = - \cos (\pi / m (s, s ^ {\prime}))
\]

并且诸 \(e_{s}\) 构成 E 的一个基。于是 W 在 E 上的自然表示可与第 3 号的表示 \(\sigma\) 等同。

